% Chapitre « Modèles » (partie 4, lot 2, jalon J6) : une section par modèle, relue avec son code.
% Sources : ADR-0009, ADR-0012, ADR-0037 ; code de src/foot_predictor/modeling/ (design.py,
% diagnostics.py, distributions.py, models/). Les chiffres cités viennent des rapports
% reports/experiments/*.json (jeu ds-2026-09-30-ba2b91f7, 4 plis, 7 156 matchs).

\chapter{Modèles et évaluation}
\label{chap:modeles}

Ce chapitre présente les modèles du MVP dans l'ordre du rapport de cadrage (I.4) : chaque étape
répond à une question. Tous partagent la même interface (apprendre sur des lignes \((m, e)\),
prédire des matchs) et la même sortie : \(\lambda_{\mathrm{dom}}\), \(\lambda_{\mathrm{ext}}\) et la
loi du total \(T = Y_{\mathrm{dom}} + Y_{\mathrm{ext}}\), repliée sur \(\{0, \dots, 9, \ge 10\}\).
Les variables sont celles des groupes G0 (contexte : domicile, championnat, huis clos) et G1 (Elo),
standardisées sur l'apprentissage du pli seulement.

% ---------------------------------------------------------------------------------------------
\section{Protocole d'évaluation et métriques}
\label{sec:protocole}

\paragraph{Plis.} Quatre plis à fenêtre croissante (ADR-0012, ADR-0037) : pour la saison de test
\(s \in \{2021, 2022, 2023, 2024\}\), apprentissage sur les lignes \((m, e)\) des saisons 2015-16 à
\(s - 1\), évaluation sur les matchs de championnat du top 5 de la saison \(s\). Les hyperparamètres se
choisissent par validation interne (apprentissage jusqu'à \(s - 2\), mesure sur \(s - 1\)), puis le
modèle est réajusté sur tout l'apprentissage du pli. Standardisation et autres ajustements se font
dans le pli ; les comparaisons portent sur les mêmes matchs (intersection).

\paragraph{Loi du total et métriques.} Chaque modèle prédit
\(\hat p_k = \hat P(T = k)\) pour \(k \in \{0, \dots, 9\}\) et \(\hat p_{10} = \hat P(T \ge 10)\). Pour un
match de total observé \(t\), avec \(c(t) = \min(t, 10)\) :
\[
  \text{log-loss} = -\log \hat p_{c(t)},
  \qquad
  \text{RPS} = \frac{1}{K - 1}\sum_{k=0}^{K-2}\Big(\hat F(k) - \mathbf{1}\{t \le k\}\Big)^2
  \quad (K = 8 \text{ catégories } 0, \dots, 6, \ge 7),
\]
\[
  \text{Brier} = \big(\hat P(T > 2{,}5) - \mathbf{1}\{t > 2{,}5\}\big)^2 .
\]
Le log-loss est la métrique principale : c'est une règle de score strictement propre, liée à la
vraisemblance, sensible à toute la loi. Une probabilité nulle sur le total observé donne une perte
infinie ; elle est refusée, jamais remplacée par un plancher.

\paragraph{Calibration.} Par décile de \(\hat p = \hat P(T \ge 3)\) (classes d'effectifs égaux), on
compare la moyenne prédite à la fréquence observée ; on estime aussi \(\operatorname{logit} P(Y = 1)
= a + b \operatorname{logit}(\hat p)\) : un modèle calibré a \(a = 0\), \(b = 1\) ; \(b < 1\) signale des
prédictions trop extrêmes. Le \emph{PIT randomisé} d'une loi discrète, \(\hat F(t - 1) + U\, \hat
p_{c(t)}\) avec \(U\) uniforme sur \([0, 1]\), est uniforme si la loi est juste.

\paragraph{Intervalle et couverture.} L'intervalle \([q_{10} ; q_{90}]\) de la loi discrète ne couvre pas
exactement 80 \% ; on publie donc sa couverture annoncée \(\hat P(T \in [q_{10} ; q_{90}])\) et l'on
compare la couverture observée à la moyenne des couvertures annoncées (ADR-0009).

\paragraph{Inférence sur les écarts.} Pour deux modèles \(A\) et \(B\), \(d_m = \ell_A(m) - \ell_B(m)\)
match par match ; \(\bar d > 0\) : \(B\) est meilleur. Les matchs d'une même journée d'un championnat
étant corrélés, l'intervalle à 95 \% vient d'un \emph{bootstrap par blocs} : on tire avec remise les
blocs (championnat, saison, journée) de chaque pli, \(10\,000\) fois, et l'on prend les percentiles
2,5 et 97,5 de \(\sum_b S_b^* / \sum_b n_b^*\) (sommes et effectifs des blocs tirés). La statistique
de Diebold-Mariano \(\bar d / \sqrt{\hat V}\), avec \(\hat V = \sum_b (S_b - n_b \bar d)^2 / N^2\), complète.

\paragraph{Règle de décision (ADR-0037, écrite avant toute évaluation).} Un modèle ou un groupe de
variables plus complexe \(B\) remplace \(A\) si (i) le gain moyen poolé des 4 plis est positif et son
intervalle à 95 \% exclut 0 ; (ii) le gain est positif dans au moins 3 plis sur 4 ; (iii) la
calibration poolée ne se dégrade pas (écart absolu moyen par décile de \(P(T \ge 3)\) : \(+0{,}5\)
point au plus ; pente dans \([0{,}9 ; 1{,}1]\)). Sinon, on garde le plus simple.

% ---------------------------------------------------------------------------------------------
\section{M1 : la régression linéaire sur le total}
\label{sec:m1}

\paragraph{Modèle.} Pour un match \(m\) de variables \(x_m\) (constante, indicatrices de
championnat, huis clos, Elo des deux équipes),
\[
  T_m = x_m^\top \beta + \varepsilon_m,
  \qquad
  \hat\beta = (X^\top X)^{-1} X^\top T .
\]
\paragraph{Hypothèses de Gauss-Markov.} (1) linéarité ; (2) exogénéité, \(\mathbb{E}[\varepsilon
\mid X] = 0\) ; (3) homoscédasticité, \(\operatorname{Var}(\varepsilon_m \mid X) = \sigma^2\) ;
(4) erreurs non corrélées ; (5) \(X\) de plein rang. Sous (1) à (5), \(\hat\beta\) est le meilleur
estimateur linéaire sans biais, de variance \(\sigma^2 (X^\top X)^{-1}\) ; \(\hat\sigma^2 = \lVert T -
X\hat\beta\rVert^2 / (n - p)\). Le coefficient de détermination est \(R^2 = 1 - \sum_m \hat\varepsilon_m^2
/ \sum_m (T_m - \bar T)^2\).

\paragraph{Intervalles.} L'intervalle de confiance de la moyenne \(x^\top\beta\) a pour écart-type
\(\mathrm{se}(x^\top\hat\beta) = \hat\sigma \sqrt{x^\top (X^\top X)^{-1} x}\) ; celui de
\emph{prédiction} d'un nouveau total ajoute le bruit :
\(s(x) = \sqrt{\hat\sigma^2 + \mathrm{se}(x^\top\hat\beta)^2}\).

\paragraph{Loi prédictive discrétisée.} Avec \(\mu = x^\top\hat\beta\) et \(s = s(x)\),
\[
  P(T = 0) = \Phi\!\Big(\frac{\tfrac12 - \mu}{s}\Big),\quad
  P(T = k) = \Phi\!\Big(\frac{k + \tfrac12 - \mu}{s}\Big) - \Phi\!\Big(\frac{k - \tfrac12 - \mu}{s}\Big)\ (1 \le k \le 9),\quad
  P(T \ge 10) = 1 - \Phi\!\Big(\frac{9{,}5 - \mu}{s}\Big).
\]
La masse sous \(-\tfrac12\), celle des totaux \og négatifs \fg{}, est rabattue sur 0.

\paragraph{Pourquoi M1 ne suffit pas (chiffres des 4 plis).} Le \(R^2\) d'apprentissage vaut
\(2{,}1\,\%\) : le total est très bruité (rapport C.2). Aucune moyenne prédite n'est négative (la plus
petite vaut \(2{,}13\)), mais la loi normale met \(2{,}5\,\%\) de sa masse sur des totaux négatifs, ce
qui gonfle \(P(T = 0)\) (\(8{,}7\,\%\) prédit contre \(6{,}4\,\%\) observé) et creuse \(P(T = 1)\). La
variance des résidus croît avec la valeur ajustée (de \(2{,}48\) à \(3{,}13\) du premier au dernier
quintile ; test de Breusch-Pagan, \(p < 10^{-5}\)) : l'hypothèse (3) est fausse, comme on l'attend
d'un comptage dont la variance suit la moyenne. L'intervalle normal est symétrique alors que la loi
d'un comptage est asymétrique à droite. Résultat : M1 est \emph{moins bon} que la référence B1 sur
le log-loss du total (écart \(-0{,}0108\), intervalle à 95 \% \([-0{,}0166 ; -0{,}0049]\)).

% ---------------------------------------------------------------------------------------------
\section{M2 : le modèle de Poisson sur le total}
\label{sec:m2}

\paragraph{Modèle.} \(T_m \sim \mathcal{P}(\mu_m)\), avec le lien logarithme
\(\log \mu_m = x_m^\top\beta\) : la moyenne est toujours positive et les effets sont
multiplicatifs. La log-vraisemblance
\[
  \ell(\beta) = \sum_m \big[T_m\, x_m^\top\beta - e^{x_m^\top\beta} - \log T_m!\big]
\]
est concave ; son maximum annule le score \(\sum_m (T_m - \mu_m)\, x_m = 0\) et s'obtient par
moindres carrés repondérés itérés (IRLS). Conséquence du score : avec une constante et des
indicatrices de championnat, la somme des \(\hat\mu_m\) de chaque championnat égale la somme des
totaux observés (propriété vérifiée par les tests).

\paragraph{Déviance et dispersion.} La déviance \(D = 2 \sum_m [T_m \log(T_m/\hat\mu_m) - (T_m -
\hat\mu_m)]\) mesure l'écart au modèle saturé. Sous Poisson, \(\operatorname{Var}(T \mid x) =
\mathbb{E}(T \mid x)\). On le vérifie par
\[
  \hat\varphi = \frac{1}{n - p} \sum_m \frac{(T_m - \hat\mu_m)^2}{\hat\mu_m}
  \quad\text{et le test de Cameron-Trivedi :}\quad
  \frac{(T_m - \hat\mu_m)^2 - T_m}{\hat\mu_m} = \alpha\, \hat\mu_m + \eta_m ,
\]
régression sans constante dont l'alternative est \(\operatorname{Var}(T) = \mu + \alpha\mu^2\)
(binomiale négative). Résultat : \(\hat\varphi = 0{,}99\) dans chaque pli, \(\hat\alpha \approx
-0{,}004\) (\(t \approx -1\)) : aucune surdispersion du total une fois le championnat et l'Elo pris
en compte.

\paragraph{Ce que M2 apporte.} Une vraie loi de comptage, positive et asymétrique : M2 bat M1 de
\(0{,}0162\) (\([0{,}0115 ; 0{,}0208]\)) et bat B1 de \(0{,}0054\) (\([0{,}0018 ; 0{,}0090]\)),
avec une pente de calibration poolée de \(0{,}97\). Ce qu'il rate : le total mélange deux équipes ;
\(\log(\lambda_{\mathrm{dom}} + \lambda_{\mathrm{ext}})\) n'est pas linéaire dans les forces des
équipes, et l'on perd \og qui marque \fg{} (section suivante, M3).

% ---------------------------------------------------------------------------------------------
\section{M3 : le modèle de Poisson par équipe}
\label{sec:m3}

\paragraph{Modèle.} Chaque match donne deux lignes \((m, e)\), une par équipe. Avec \(x_{m,e}\) les
variables de l'équipe \emph{et} de son adversaire (constante, championnat, domicile, huis clos et son
interaction avec le domicile, Elo de l'équipe et Elo de l'adversaire),
\[
  Y_{m,e} \sim \mathcal{P}(\lambda_{m,e}),
  \qquad
  \log \lambda_{m,e} = x_{m,e}^\top \beta .
\]
Le lien log rend multiplicative la rencontre de l'attaque d'une équipe et de la défense de
l'autre : \(\lambda_{m,e} = e^{\beta_0}\, e^{\beta_{\mathrm{dom}} d_{m,e}}\, e^{\beta_{\mathrm{Elo}} z_e}\,
e^{\beta_{\mathrm{adv}} z_{\mathrm{adv}(e)}} \cdots\).

\paragraph{Loi du total.} Si \(Y_{\mathrm{dom}}\) et \(Y_{\mathrm{ext}}\) sont indépendants sachant
les variables,
\[
  P(T = k) = \sum_{a=0}^{k} P(Y_{\mathrm{dom}} = a)\, P(Y_{\mathrm{ext}} = k - a),
  \qquad\text{soit}\qquad T \sim \mathcal{P}(\lambda_{\mathrm{dom}} + \lambda_{\mathrm{ext}}),
\]
et la loi jointe \(P(a, b) = P(Y_{\mathrm{dom}} = a) P(Y_{\mathrm{ext}} = b)\) donne le score exact
et le 1N2.

\paragraph{Erreurs standard groupées.} Les deux lignes d'un match ne sont pas indépendantes. On
estime \(\operatorname{Var}(\hat\beta)\) par le sandwich par grappes
\(\hat V = A^{-1} \big(\sum_{m} s_m s_m^\top\big) A^{-1}\), où \(A\) est l'information de Fisher et
\(s_m = \sum_{e \in m} (y_{m,e} - \hat\lambda_{m,e}) x_{m,e}\) le score du match. Sur nos données,
elles valent de \(0{,}97\) à \(1{,}00\) fois les erreurs \og naïves \fg{} : sachant les variables,
les deux lignes d'un match sont très légèrement corrélées négativement.

\paragraph{Ce que M3 apporte.} Sur les 4 plis, M3 bat M2 de \(0{,}0033\) (\([0{,}0014 ;
0{,}0052]\), 4 plis sur 4) et B1 de \(0{,}0087\) (\([0{,}0048 ; 0{,}0126]\)) ; pente de calibration
poolée \(1{,}01\). Décomposer par équipe améliore la loi du total. Coefficients (pli 2024-25) :
\(\hat\beta_{\mathrm{dom}} = 0{,}25\) (une équipe à domicile marque \(e^{0{,}25} \approx 1{,}28\) fois
plus) ; à huis clos, l'interaction vaut \(-0{,}15\) : l'avantage tombe à \(e^{0{,}10} \approx 1{,}11\).

% ---------------------------------------------------------------------------------------------
\section{M4 : la dispersion, et pourquoi la binomiale négative n'est pas retenue}
\label{sec:m4}

\paragraph{Modèle.} La binomiale négative NB2 garde la moyenne \(\mu\) et ajoute un paramètre
\(\alpha \ge 0\) :
\[
  P(Y = y) = \frac{\Gamma(y + \alpha^{-1})}{\Gamma(\alpha^{-1})\, y!}
  \Big(\frac{\alpha^{-1}}{\alpha^{-1} + \mu}\Big)^{\alpha^{-1}}
  \Big(\frac{\mu}{\alpha^{-1} + \mu}\Big)^{y},
  \qquad
  \operatorname{Var}(Y) = \mu + \alpha \mu^2 .
\]
Elle se justifie en cas de \emph{surdispersion}. On la teste par Cameron-Trivedi (section
\ref{sec:m2}) et par le rapport de vraisemblance \(LR = 2(\ell_{NB} - \ell_{\mathcal{P}})\). Sous
\(H_0 : \alpha = 0\), le paramètre est au bord de son domaine : \(LR\) suit le mélange
\(\tfrac12 \chi^2_0 + \tfrac12 \chi^2_1\), d'où \(p = \tfrac12 P(\chi^2_1 > LR)\).

\paragraph{Constat.} Sur l'apprentissage de chacun des 4 plis : \(\hat\varphi\) de \(0{,}976\) à
\(0{,}980\), \(\hat\alpha_{CT}\) de \(-0{,}017\) à \(-0{,}014\) (\(p\) de \(0{,}001\) à \(0{,}017\)) :
les buts d'une équipe sont légèrement \emph{sous}-dispersés une fois les variables connues. La NB2
estime \(\alpha\) au bord (\(\approx 10^{-7}\)), \(LR \approx 0\), \(p = 0{,}5\). La binomiale négative
ne peut pas représenter une sous-dispersion : elle se confond avec Poisson. Elle n'est donc pas
évaluée sur les plis (règle de la sous-étape 4.8), et M3 reste la structure de référence.

% ---------------------------------------------------------------------------------------------
\section{M5 : la dépendance entre les deux équipes (Dixon-Coles)}
\label{sec:m5}

\paragraph{Modèle.} M3 suppose \(Y_{\mathrm{dom}}\) et \(Y_{\mathrm{ext}}\) indépendants sachant
les variables. Dixon et Coles (1997) corrigent les quatre scores les plus bas :
\[
  P(a, b) = \tau_{\lambda,\mu}(a, b)\; \frac{e^{-\lambda}\lambda^{a}}{a!}\; \frac{e^{-\mu}\mu^{b}}{b!},
  \qquad
  \tau(a, b) =
  \begin{cases}
    1 - \lambda\mu\rho & (a, b) = (0, 0)\\
    1 + \lambda\rho & (a, b) = (0, 1)\\
    1 + \mu\rho & (a, b) = (1, 0)\\
    1 - \rho & (a, b) = (1, 1)\\
    1 & \text{sinon,}
  \end{cases}
\]
avec \(\lambda = \lambda_{\mathrm{dom}}\), \(\mu = \lambda_{\mathrm{ext}}\). Les corrections se
compensent (\(-\lambda\mu\rho\, p_{00} + \lambda\rho\, p_{01} + \mu\rho\, p_{10} - \rho\, p_{11} = 0\)
pour des marges de Poisson) : \(P\) reste une loi. La loi du total vient de la loi jointe,
\(P(T = k) = \sum_{a + b = k} P(a, b)\). La contrainte \(\tau > 0\) impose
\(\max(-1/\lambda, -1/\mu) < \rho < \min(1/(\lambda\mu), 1)\) pour chaque match.

\paragraph{Estimation.} En deux temps, dans le pli : les \(\lambda\) de M3 par maximum de
vraisemblance, puis \(\hat\rho = \arg\max_\rho \sum_m \log \tau(y_{m,\mathrm{dom}}, y_{m,\mathrm{ext}};
\hat\lambda_m, \hat\mu_m, \rho)\).

\paragraph{Résultat.} \(\hat\rho\) vaut de \(-0{,}061\) à \(-0{,}055\) selon le pli : les scores 0-0
et 1-1 sont un peu plus fréquents que sous l'indépendance. Sur la loi du total (4 plis poolés),
\(P(T = 1)\) passe de \(17{,}6\,\%\) à \(16{,}3\,\%\) (observé : \(16{,}3\,\%\)), mais \(P(T = 0)\) de
\(6{,}5\,\%\) à \(7{,}1\,\%\) (observé : \(6{,}3\,\%\)) : la correction répare \(T = 1\) et dépasse sur
\(T = 0\). Le log-loss du total ne bouge pas (\(+0{,}00002\), \([-0{,}0009 ; +0{,}0010]\)) : la
corrélation change la forme de la loi sur les petites valeurs, sans gain global. Par la règle de
décision, M3 (plus simple) est conservé.

% ---------------------------------------------------------------------------------------------
\section{M6 : plus de variables, sous régularisation}
\label{sec:m6}

\paragraph{Variables.} Aux groupes G0 et G1 s'ajoutent G2 (logarithmes des moyennes glissantes de
buts et d'\(\mathrm{xg\_proxy}\), pour et contre, de l'équipe et de l'adversaire, à la demi-vie
\(h\)) et G3 (repos borné à 14 jours, matchs des 14 derniers jours, match européen dans les 4 jours,
pour les deux équipes). Toutes les variables continues sont standardisées sur l'apprentissage du
pli : sans cela, la pénalité frapperait différemment des variables d'unités différentes.

\paragraph{Ridge.} Avec \(D(y, \beta) = 2\sum_i [y_i \log(y_i/\lambda_i) - (y_i - \lambda_i)]\) la
déviance de Poisson et \(\beta_0\) la constante, non pénalisée,
\[
  \hat\beta(\alpha) = \operatorname*{arg\,min}_{\beta}\ \frac{1}{2n} D(y, \beta)
  + \frac{\alpha}{2} \sum_{j \ge 1} \beta_j^2 .
\]
\paragraph{Élastique net.} La pénalité devient
\(\alpha \big[\omega \sum_j |\beta_j| + \tfrac{1 - \omega}{2} \sum_j \beta_j^2\big]\), avec
\(\omega = 0{,}5\) ; la partie \(L_1\) peut annuler des coefficients (sélection de variables).

\paragraph{Compromis biais-variance.} Quand \(\alpha\) croît, les coefficients rétrécissent vers 0
(variance plus faible, biais plus fort) ; \(\alpha \to 0\) redonne le Poisson non régularisé
(propriété vérifiée par les tests), \(\alpha \to \infty\) une prédiction constante. \(\alpha\) et
\(h \in \{60, 120, 240\}\) jours se choisissent par validation interne : apprentissage jusqu'à
l'avant-dernière saison du pli, mesure du log-loss du total sur la dernière, puis réajustement.

\paragraph{Non-linéarité.} Avant d'ajouter des splines, on compare, par décile de l'écart d'Elo,
les buts observés \(O\) aux buts attendus \(E = \sum \hat\lambda\) de M3, sur 2015-16 à 2020-21
seulement (jamais une saison de test) :
\(z = (O - E)/\sqrt{E}\) et \(\sum z^2 \approx \chi^2_9\). On obtient \(\sum z^2 = 9{,}86\)
(\(p = 0{,}36\)) et des rapports \(O/E\) entre \(0{,}97\) et \(1{,}03\), sans forme : la relation
log-linéaire à l'Elo suffit, et aucune spline n'est ajoutée.

\paragraph{Résultat.} Avec les mêmes variables G0 à G3, le ridge ne bat pas le Poisson non régularisé
(\(+0{,}0002\), \([-0{,}0004 ; +0{,}0008]\)), l'élastique net non plus (\(+0{,}0001\),
\([-0{,}0001 ; +0{,}0004]\)) : avec environ \(33\,000\) lignes pour une trentaine de coefficients, la
variance d'estimation est déjà faible. En revanche, les variables G2 et G3 font gagner \(0{,}0091\)
(\([0{,}0056 ; 0{,}0126]\)) au Poisson par équipe. La structure retenue reste M3 (non régularisé) ; les
ablations du lot 3 mesurent la part de chaque groupe.

% ---------------------------------------------------------------------------------------------
\section{Ablations et modèle retenu}
\label{sec:retenu}

\paragraph{Ablations.} Sur la structure M3, les groupes s'ajoutent un à un et chaque ajout est jugé
par la règle : G0 (contexte) gagne \(0{,}0039\) sur B0 ; G1 (Elo) \(0{,}0094\)
(\([0{,}0060 ; 0{,}0128]\)) de plus ; G2 (moyennes glissantes) \(0{,}0092\) (\([0{,}0057 ; 0{,}0127]\)) de
plus ; G3 (calendrier) \(-0{,}0001\) (\([-0{,}0004 ; 0{,}0002]\)) : il n'est pas retenu. Apprendre aussi
sur les deuxièmes divisions n'apporte rien de significatif (\(+0{,}0003\), \([-0{,}0003 ; 0{,}0010]\)).

\paragraph{Modèle MVP, horizon H1 (ADR-0039).} Pour une ligne \((m, e)\),
\[
  \log \lambda_{m,e} = \beta_0 + \beta_{\mathrm{ch}(m)} + \beta_{\mathrm{dom}}\, d_{m,e}
  + \beta_{\mathrm{hc}}\, h_m + \beta_{\mathrm{dom \times hc}}\, d_{m,e} h_m
  + \beta_{1}\, \tilde R_e + \beta_{2}\, \tilde R_{\mathrm{adv}(e)}
  + \sum_{j=1}^{8} \gamma_j\, \widetilde{\log g_j},
\]
où \(\tilde R\) est l'Elo standardisé, et \(\widetilde{\log g_j}\) les logarithmes standardisés des huit
moyennes glissantes (buts et \(\mathrm{xg\_proxy}\), pour et contre, de l'équipe et de l'adversaire) à
la demi-vie choisie par validation interne. \(Y_{m,e} \sim \mathcal{P}(\lambda_{m,e})\), estimé par
maximum de vraisemblance sur le top 5, et \(T \sim \mathcal{P}(\lambda_{\mathrm{dom}} +
\lambda_{\mathrm{ext}})\). Sur les 4 plis, il bat B1 de \(0{,}0179\) (\([0{,}0133 ; 0{,}0226]\)), avec une
pente de calibration poolée de \(0{,}95\) et une couverture de \([q_{10} ; q_{90}]\) à \(\pm 1{,}1\) point
de la couverture annoncée. Le même modèle sert en rejeu et en live (G3 non retenu). Le test scellé,
sur les matchs joués depuis le 1\textsuperscript{er} juillet 2025, dira s'il tient ses promesses ; son
résultat sera publié tel quel.
